\section{Theorem dependencies}\label{app:dependencies}

The following table records the logical inputs to the principal statements.
``Exact algebra'' means an identity proved in the indicated polynomial or
localized coordinate ring; bounded computations are not dependencies.

\small
\begin{longtable}{>{\raggedright\arraybackslash}p{0.20\textwidth}
                  >{\raggedright\arraybackslash}p{0.34\textwidth}
                  >{\raggedright\arraybackslash}p{0.34\textwidth}}
\toprule
Result & Internal dependencies & External input \\
\midrule
\endfirsthead
\toprule
Result & Internal dependencies & External input \\
\midrule
\endhead
Theorem~\ref{thm:geometric-counterexample} &
Proposition~\ref{prop:geometric-affine-source}; chord identities
\eqref{eq:auxiliary-chords} &
Symmetric powers of $\P^1$; finite flat marked-point cover; triviality of line
bundles and additive torsors on affine space \\
\addlinespace
Theorem~\ref{thm:counterexample} &
Coordinate identity \eqref{eq:trc-target}, chain rule, exact substitution &
Jacobian criterion for etaleness over a characteristic-zero field \\
\addlinespace
Theorem~\ref{thm:marked-root} &
Proposition~\ref{prop:cubic-marking}; Lemmas~\ref{lem:affine-mutual} and
\ref{lem:projective-mutual}; overlap \eqref{eq:chart-overlap} &
None; this is an explicit two-chart coordinate certificate \\
\addlinespace
Proposition~\ref{prop:fibers-image} &
Theorem~\ref{thm:marked-root}; multiplicity classification of a binary cubic;
coefficient comparison for a cube & None beyond elementary algebra \\
\addlinespace
Proposition~\ref{prop:nonproperness} &
Both cubic reconstruction charts; discriminant identity &
Complex implicit-function theorem and the sequential criterion for
nonproperness \\
\addlinespace
Proposition~\ref{prop:weighted-polynomiality} &
Endpoint conditions \eqref{eq:admissible-H}; divisibility of $p,q$; divided
differences & None \\
\addlinespace
Proposition~\ref{prop:weighted-jacobian} &
Theorem~\ref{thm:tangent-core} and its four coordinate-change determinants &
Etale morphisms of finite type are quasi-finite \\
\addlinespace
Proposition~\ref{prop:weighted-Cnonzero} &
Incidence equation \eqref{eq:weighted-pencil}; explicit reconstruction
\eqref{eq:weighted-reconstruction} & None \\
\addlinespace
Theorem~\ref{thm:weighted-global} &
Propositions~\ref{prop:weighted-polynomiality},
\ref{prop:weighted-jacobian}, and \ref{prop:weighted-Cnonzero} &
Finiteness and universal property of normalization; Zariski's Main Theorem;
normality criterion via height-one valuations \\
\addlinespace
Proposition~\ref{prop:pencil-irreducible} &
Linearity in $t$ and coefficient comparison & Gauss's lemma \\
\addlinespace
Proposition~\ref{prop:universal-disc} &
Repeated-root equations and coprime pole orders $n-1,n$ &
Standard facts on finite maps of normal projective curves \\
\addlinespace
Theorem~\ref{thm:symmetric-monodromy} &
Lemma~\ref{lem:generic-morse-slice}; restriction to a vertical line &
Classical Morse-polynomial monodromy theorem \\
\bottomrule
\end{longtable}
\normalsize
